\SourceLeaf{226}%
\bigskip\Needspace{9\baselineskip}
\pdfbookmark[1]{Pallas Olbersiana. August 1802}{pallas180208}
\begin{center}
PALLAS OLBERSIANA.\par
\medskip\hrule\medskip
Monatliche Correspondenz. Vol. VI. pp.\ 187—196. 1802 August.\par
\medskip\hrule
\end{center}\medskip
\begingroup\emergencystretch=2em
—\quad —\quad —\quad We have also received confirmation from \textit{Vilna}, \textit{Prague}, \textit{Kremsmünster} and \textit{Padua} that \textit{Pallas} was not observed at those places. According to a report from Dr.\ \textsc{Gauss}, Dr.\ \textsc{Maskelyne} had still not given up observing this planet on 20 July. He believes that on 18 July he still managed to observe \textit{Pallas} as a star of the 11th magnitude, though he was not quite sure whether it was the planet or a fixed star. Calculation or the next observation will show which it was.\par\endgroup

Dr.\ \textsc{Gauss}’s elements III for the orbit of \textit{Pallas} still agree well with the observed declinations. Only the right ascensions now differ by a few seconds more. Dr.\ \textsc{Gauss} compared the observations that \textsc{Cesaris} made in Milan with an excellent eight-foot mural quadrant by \textsc{Ramsden}. We published them on p.\ 73 of the previous issue. He found the following agreement.
\begin{center}
\textit{Comparison of the Milan meridian observations of Pallas}\par
\textit{with }\textsc{Gauss}\textit{’s elements III.}\par\medskip
{\small\setlength{\tabcolsep}{4pt}
\begin{tabular}{r r|r|r|r|r|r}
\multicolumn{2}{c|}{1802}&\shortstack{Mean time\\in Milan}&\multicolumn{2}{c|}{Calculated}&\multicolumn{2}{c}{\shortstack{Difference\\R—B}}\\
\multicolumn{2}{c|}{}&&\shortstack{Right\\ascension}&\shortstack{Northern\\declination}&in R.A.&in decl.\\\hline
May&4&$9^{\mathrm{u}}\ 15^{\mathrm{m}}\ 43^{\mathrm{s}}1$&$180^0\ 56'\ 15''7$&$20^0\quad2'\ 44''1$&$+15''7$&$-5''9$\\
&5&$9\ 11\ 45.3$&$180\ 55\ 56.8$&$20\quad8\ 38.1$&$+19.8$&$+5.1$\\
&6&$9\quad7\ 49.3$&$180\ 56\quad0.3$&$20\ 14\ 10.4$&$+16.3$&$+6.4$\\
&7&$9\quad3\ 55.6$&$180\ 56\ 24.8$&$20\ 19\ 22.1$&$+10.8$&$+2.1$\\
&8&$9\quad0\quad3.2$&$180\ 57\ 13.3$&$20\ 24\ 12.4$&$+8.3$&$+0.4$\\
&9&$8\ 56\ 12.1$&$180\ 58\ 21.9$&$20\ 28\ 43.1$&$+6.9$&$+0.1$\\
&10&$8\ 52\ 22.4$&$180\ 59\ 52.8$&$20\ 32\ 53.0$&$+7.8$&$-2.0$\\
&11&$8\ 48\ 34.3$&$181\quad1\ 45.2$&$20\ 36\ 43.4$&$+9.2$&$-3.6$\\
&12&$8\ 44\ 47.0$&$181\quad3\ 58.8$&$20\ 40\ 14{,}6$&$+7.8$&$-1.4$\\
&17&$8\ 26\ 12.9$&$181\ 20\ 17.1$&$20\ 53\ 18.0$&$+11.1$&$-4.0$\\
&18&$8\ 22\ 33.1$&$181\ 24\ 33.5$&$20\ 55\quad3.6$&$+15.5$&$+0.6$\\
&19&$8\ 18\ 55.8$&$181\ 29\quad9.4$&$20\ 56\ 30.8$&$+16.4$&$-0.2$\\
&20&$8\ 15\ 19.6$&$181\ 34\quad4.8$&$20\ 57\ 43.0$&$+14.8$&$-2.0$\\
&21&$8\ 11\ 44.6$&$181\ 39\ 18.6$&$20\ 58\ 39.5$&$+10.6$&$-1.5$\\
&22&$8\quad8\ 11.1$&$181\ 44\ 52.9$&$20\ 59\ 20.6$&$+6.9$&$-1.4$
\end{tabular}}
\end{center}

\Needspace{22\baselineskip}
Dr.\ \textsc{Olbers}’s observations show the same agreement. Their consistency with one another also shows the exceptional value of these observations, even though they were made with only a \textit{ring micrometer}. Anyone who takes the trouble to examine them as closely as we have will find that they are in no way inferior to those made with equatorial sectors, that they stand alongside them and that they are sometimes even preferable. This, of course, says more for the skill of the excellent observer than for the quality of the instrument, as our readers will see from the following comparison.
\newcommand{\PallasSymbol}{\text{\mdseries\AstroSymbols ⚴}}
\SourceLeaf{227}%
\Needspace{15\baselineskip}
\begin{center}
\textit{Comparison of Dr.\ }\textsc{Olbers}\textit{’s observations of Pallas}\par
\textit{with }\textsc{Gauss}\textit{’s elements III.}\par\medskip
{\small\setlength{\tabcolsep}{4pt}
\begin{tabular}{r r|r|r|r|r|r}
\multicolumn{2}{c|}{1802}&\shortstack{Mean time\\in Bremen}&\multicolumn{2}{c|}{Calculated}&\multicolumn{2}{c}{\shortstack{Difference\\R—B}}\\
\multicolumn{2}{c|}{}&&\shortstack{Right ascension\\of $\PallasSymbol$}&\shortstack{Right ascension\\of $\PallasSymbol$}&in R.A.&in decl.\\\hline
June&19&$11^{\mathrm{u}}\ 13^{\mathrm{m}}\ 55^{\mathrm{s}}$&$186^0\ 16'\ 19''$&$19^0\ 54'\ 17''$&$-6''$&$+17''$\\
&20&$11\quad7\ 51$&$186\ 29\ 12$&$19\ 49\ 39$&$-22$&$-1$\\
&21&$11\quad8\ 54$&$186\ 42\ 39$&$19\ 44\ 55$&$-20$&$-3$\\
&26&$11\ 18\ 34$&$187\ 52\quad3$&$19\ 19\ 31$&$-16$&$+8$\\
July&4&$11\ 28\ 14$&$189\ 52\ 12$&$18\ 33\ 39$&$-26$&$-5$\\
&8&$11\quad5\ 16$&$190\ 55\ 44$&$18\quad8\ 50$&$-25$&$-8$\\
&8&$11\ 13\ 20$&$190\ 55\ 50$&$18\quad8\ 48$&$-23$&$-10$\\
&9&$11\ 53\ 47$&$191\ 12\ 36$&$18\quad2\ 35$&$-28$&$+1$
\end{tabular}}
\end{center}

All these comparisons and the excellent agreement of the observations with \textsc{Gauss}’s ellipse (which this skilful calculator will bring still closer to the truth) give us grounds to hope with confidence that we will keep this new planet even if its faintness makes it impossible to find in 1803 and 1804. Even now, however, we have sound reason to hope that \textit{Pallas} will not escape our keen eyes next year. Dr.\ \textsc{Olbers}, who last saw the planet on 9 July, is himself of this opinion and writes, “However hard the planet is to see now, I am convinced that you will see it next year with your transit instrument. At opposition it will then be a little closer to the Earth than it was on 9 July. It is farther from the Sun and is more faintly illuminated by it, but its greater height above the horizon and its greater distance from the twilight more than make up for this.”

\Needspace{23\baselineskip}
Dr.\ \textsc{Gauss} is also very hopeful of finding \textit{Pallas} again next year. He finds its distance from the Earth on 28 Jan.\ 1803 (about one day before opposition) to be $2.56$. From Dr.\ \textsc{Herschel}’s own measurement he finds a somewhat larger diameter for this planet. On 22 April, Dr.\ \textsc{Herschel} gives the apparent diameter as $=0''17$, based on a fairly good measurement. From this Dr.\ \textsc{Gauss} calculates the true diameter as $26\frac{1}{2}$ German miles (the distance from the Earth being $=1.562$). In his latest letter he expresses surprise at the enormous difference between the results of Dr.\ \textsc{Herschel}’s and Dr.\ \textsc{Schröter}’s measurements of these diameters, even though they used the same method. “I am very curious,” he writes, “what magnifications Dr.\ \textsc{Herschel} used. Surely 500 times magnification would still not make an apparent diameter of $0''17$ show up as a disk?” At 300 times magnification I, at least, could see no trace of a disk on either this planet of \textsc{Olbers} or that of \textsc{Piazzi}. I asked Dr.\ \textsc{Olbers} about it and received this reply. “I too have never been able, even with the highest magnifications of my \textsc{Dollond} achromat, to see a difference between these planets and fixed stars of the same apparent magnitude.”

In the previous issue, p.\ 85 [p.\ 219 of this volume], we published an ephemeris for \textit{Pallas} by \textsc{Gauss}. Here is its continuation. It appears too late to help with the observations, but it will still help prevent any stars mistaken for the planet from being reported as \textit{Pallas}. [The ephemeris is arranged in the same way as the one printed on p.\ 219, except that it omits the distance of Pallas from the Sun and runs through 28 August 1802. The ephemerides included so far in this edition of \textsc{Gauss}’s works will suffice as examples of those he calculated. The rest may be left out to save space.]
